% Moved from sections/library.tex (PoPETS-2027/main.tex lines 255-678) in the
% S&P version: the Okamoto, Neq and linear-relation protocols and the
% listings that the compressed Section 3 refers to.
\section{Further Protocols of the Library}\label{app:eqneq}
This appendix describes the Okamoto, Neq and linear-relation protocols of
the library and collects the listings that Section~\ref{schn_proof}
refers to.

\begin{lstlisting}[frame=single, language=Coq, caption={Completeness and Soundness},
  label={comp_sound},captionpos=t, basicstyle=\ttfamily\scriptsize,
  abovecaptionskip=-\medskipamount]
Lemma schnorr_completeness (g h : G) (x : F) (R : h = g^x) :
  forall (u c : F),
  accepting_conversation g h (schnorr_protocol x g u c) = true.
Proof. (* proof terms omitted *) Qed.

Lemma simulator_completeness (g h : G) :
  forall (u c : F),
  accepting_conversation g h (schnorr_simulator g h u c) =
    true.
Proof. (* proof terms omitted *) Qed.

Lemma special_soundness_berry_gen (g h : G) :
  forall (a : G) (c1 r1 c2 r2 : F),
  c1 <> c2 ->
  accepting_conversation g h ([a]; [c1]; [r1]) = true ->
  accepting_conversation g h ([a]; [c2]; [r2]) = true ->
  exists y : F, g^y = h /\ y = ((r1 - r2) * inv (c1 - c2)).
Proof. (* proof terms omitted *) Qed.
\end{lstlisting}

\begin{lstlisting}[frame=single, language=Coq, caption={Uniform Probability Distribution},
    label={prob_def},captionpos=t, basicstyle=\ttfamily\scriptsize,
    abovecaptionskip=-\medskipamount]
(* Probability Distribution on a type A *)
Definition dist (A : Type) : Type := list (A * prob).

(* Two distributions are equal when they are
   permutations of each other *)
Definition dist_equiv {A : Type} : relation (dist A) :=
  fun xs ys => Permutation xs ys.

(* Probability Monad *)
Definition Ret {A : Type} (x : A) : dist A := [(x, one)].

Fixpoint Bind {A B : Type} (xs : dist A)
  (f : A -> dist B) : dist B :=
  match xs with
  | [] => []
  | (ax, px) :: tx =>
    (List.map
      (fun '(ut, pt) => (ut, mul_prob px pt)) (f ax)) ++
    Bind tx f
  end.

(* Uniform distribution *)
Definition uniform_with_replacement {A : Type} :
  forall (l : list A), l <> [] -> dist A.
Proof.
  intros ? H.
  remember (Pos.of_nat (List.length l)) as len.
  exact (List.map (fun x =>  (x, mk_prob 1 len)) l).
Defined.
\end{lstlisting}

\begin{lstlisting}[frame=single, language=Coq, caption={Composition of a Distribution},
label={dist_comp},captionpos=t, basicstyle=\ttfamily\scriptsize,
abovecaptionskip=-\medskipamount]
Fixpoint repeat_dist_ntimes_vector {A : Type} (d : dist A)
  (n : nat) : dist (Vector.t A n) :=
  match n with
  | 0 => Ret (@Vector.nil A)
  | S n' =>
    Bind d (fun u =>
      Bind (repeat_dist_ntimes_vector d n')
      (fun v => Ret (Vector.cons _ u _ v)))
  end.

Lemma uniform_probability_multidraw_prob {A : Type} :
  forall n lf (a : Vector.t A n) b (Hlf : lf <> []),
  In (a, b) (repeat_dist_ntimes_vector
    (uniform_with_replacement lf Hlf) n) ->
  b = mk_prob 1 (Pos.of_nat (Nat.pow (List.length lf) n)).
Proof. (* proof terms omitted *) Qed.

Definition generalised_parallel_schnorr_distribution
  {n : nat} (lf : list F) (Hlfn : lf <> List.nil)
  (x : F) (g : G) (cs : Vector.t F n) :
  dist (@sigma_proto F G n n n) :=
  (* draw n random elements *)
  us <- repeat_dist_ntimes_vector
    (uniform_with_replacement lf Hlfn) n ;;
  Ret (construct_parallel_conversations_schnorr x g us cs).
\end{lstlisting}

\subsection{Okamoto Protocol}

Okamoto's protocol~\cite[Identification scheme 1]{okamoto1992provably} was originally introduced as a provably secure variation of Schnorr's identification protocol.
While Schnorr's protocol proves knowledge of a single exponent $x$ such that $h = g^x$, Okamoto's protocol proves knowledge of
exponents $x_1, x_2$ satisfying $h = g_1^{x_1} \cdot g_2^{x_2}$, given fixed generators $g_1,g_2$. Okamoto's protocol and its generalisation to multiple exponents and generators
(already considered in~\cite{10.1007/3-540-39118-5_13}) are commonly viewed as $\Sigma$-protocols, and are often used as building block in more advanced constructions.
We will use it, in particular, for NEQ-composition as explained in the next section.

The generalised Okamoto protocol follows the standard 3-move $\Sigma$-protocol structure. Given public generators $\mathbf{g} = (g_1, g_2, \ldots, g_n)$, public value $h$, and secret witnesses $\mathbf{x} = (x_1, x_2, \ldots, x_n)$ satisfying $h = \prod_{i=1}^{n} g_i^{x_i}$, the prover begins by selecting random values from the underlying field $\mathbf{u} = (u_1, u_2, \ldots, u_n)$ and computing announcement $t = \prod_{i=1}^{n} g_i^{u_i}$, which is sent to the verifier. Next, the verifier responds with a random challenge $c$ from the field. Finally, the prover computes responses $r_i = u_i + c \cdot x_i$ for each $i \in {1, \ldots, n}$ and sends the vector $\mathbf{r} = (r_1, r_2, \ldots, r_n)$ to the verifier for verification. The verifier accepts if and only if
$\prod_{i=1}^{n} g_i^{r_i} = t \cdot h^c$. Completeness follows from the fact that:
\[
\prod_{i=1}^{n} g_i^{r_i} = \prod_{i=1}^{n} g_i^{u_i + c \cdot x_i} =
\left(\prod_{i=1}^{n} g_i^{u_i}\right) \cdot \left(\prod_{i=1}^{n} g_i^{x_i}\right)^c = t \cdot h^c
\]

Our Rocq formalisation implements the generalised Okamoto protocol for $n \geq 2$ generators and can be accessed in the file \href{https://anonymous.4open.science/r/SigmaProtocol-EBA5/src/Crypto/Okamoto.v}{\textcolor{blue}{Okamoto.v}}. The key functions include: \texttt{\seqsplit{generalised\_okamoto\_commitment}} (computes $\prod_{i=1}^{n} g_i^{u_i}$), \texttt{\seqsplit{generalised\_okamoto\_response}} (computes responses $r_i = u_i + c \cdot x_i$), and
\texttt{\seqsplit{generalised\_okamoto\_accepting\_conversation}} (verifies the product equality condition). We also provide specialised instances for the base case of exactly two generators ($n=2$), which is heavily used in the Neq composition. The formalisation proves all three security properties: (i) \textbf{completeness}:
when both prover and verifier follow the protocol honestly with $h = \prod_{i=1}^{n} g_i^{x_i}$,
the verification check always succeeds. The proof proceeds by induction on $n$, leveraging the
vector space properties of scalar multiplication distributivity over field addition; (ii) \textbf{special soundness}:
given two accepting transcripts $(t, c_1, \mathbf{r}_1)$ and $(t, c_2, \mathbf{r}_2)$ with
the same commitment $t$ but different challenges $c_1 \neq c_2$, an extractor can compute
witnesses $\mathbf{x} = (x_1, \ldots, x_n)$ where
$x_i = (r_{1,i} - r_{2,i}) \cdot (c_1 - c_2)^{-1}$ for each $i$. The proof establishes that
these extracted values satisfy $h = \prod_{i=1}^{n} g_i^{x_i}$; and (iii) \textbf{special honest-verifier zero-knowledge}:
the real transcript distribution (involving secrets $\mathbf{x}$) is identical to the simulator
distribution (without secrets). The simulator chooses random $\mathbf{u}$ and $c$, then computes
the commitment as $t = \left(\prod_{i=1}^{n} g_i^{u_i}\right) \cdot h^{-c}$ with responses
$\mathbf{r} = \mathbf{u}$. Both distributions produce accepting transcripts with probability
$1/|\mathbb{F}|^n$ where $\mathbb{F}$ is the underlying field, proven using our uniform probability
distribution framework.

Beyond the standard $\Sigma$-protocol properties, our formalisation includes witness indistinguishability. We prove that transcripts generated with different witness
vectors $\mathbf{x}_1$ and $\mathbf{x}_2$ satisfying the same relation $h = \prod_{i=1}^{n} g_i^{x_{1,i}} =
\prod_{i=1}^{n} g_i^{x_{2,i}}$ are information-theoretically indistinguishable. The proof \texttt{\seqsplit{generalised\_okamoto\_witness\_indistinguishable}}, shown in Listing~\ref{wit_ind} in Appendix~\ref{appendix_listings}, relies on two
bijective transformation functions that map
randomness between the two witness spaces, along with the homomorphic property showing that
$\prod_{i=1}^{n} g_i^{u_i - v_i} = \left(\prod_{i=1}^{n} g_i^{u_i}\right) \cdot
\left(\prod_{i=1}^{n} g_i^{-v_i}\right)$.

\subsection{NEQ Composition}
The Neq composition formalises the statement that a prover knows witnesses
$x_1, \ldots, x_{2+n}$ for public instances $(g_i, h_i)$ such that $h_i = g_i^{x_i}$ for every
$i$, and moreover these witnesses are pairwise distinct, i.e., $x_i \neq x_j$ for all distinct pairs $i \neq j$.
In contrast to Or and And compositions, which combine existing statements by logical connectives, the Neq protocol has to encode a family of inequality constraints between witnesses.  In our development, this is done by combining an And proof of all discrete-logarithm relations with one Okamoto proof for every unordered pair of witnesses.

The protocol structure for $2+n$ statements (with $n \ge 0$, i.e., at least two statements) is inherently quadratic in complexity due to the pairwise nature of inequality constraints. For each pair of indices $(i, j)$ with $i < j$, the prover must demonstrate that $x_i \neq x_j$ without revealing the values. Our implementation works as follows: in the commitment phase, the prover generates randomness of size $(2 + n) + ((2 + n) \cdot (1 + n))$ for two components. The first $2+n$ random values $u_i$ are used for the AND component, creating commitments $t_i = g_i^{u_i}$ for each $i \in \{1,\dots,2+n\}$, thereby proving knowledge of all discrete logarithms. The remaining randomness, of size $(2+n)(1+n)$, is split into two halves of equal length and zipped together to provide two random field elements per unordered pair $(i,j)$. For each pair, the prover generates an Okamoto commitment using the combined generators $(g_i \cdot g_j,\; h_i \cdot h_j)$ and the two allocated randomness values.
The number of such commitments is $\binom{2+n}{2}$, one for each  pair $(i,j)$. The combined commitment vector, of dimension $(2+n) + \binom{2+n}{2}$, is then sent to the verifier. In the challenge phase, the verifier responds with a random challenge $c$ from the underlying field. In the response phase, the prover computes two sets of responses. The AND responses are $r_i = u_i + c \cdot x_i$ for each $i$ (as in the standard AND composition). For the Okamoto component, the prover constructs, for each pair $(i,j)$, the witnesses $w_{ij} = x_i \cdot (x_i - x_j)^{-1}$ and $w_{ji} = (x_j - x_i)^{-1}$ that encode the inequality $x_i \neq x_j$. Using these witnesses, the prover computes a two-dimensional Okamoto response; this yields $(2+n)(1+n)$ field elements. The total response vector thus has dimension $(2+n) + (2+n)(1+n) = (2+n)^2$, which is sent to the verifier.

The verifier accepts if both sub-protocol verifications succeed: (i) the AND verifier confirms that each $g_i^{r_i} = t_i \cdot h_i^c$, and (ii) for each pair $(i, j)$, the Okamoto verifier confirms that the relation $(g_i \cdot g_j)^{r_{i,1}} \cdot (h_i \cdot h_j)^{r_{i,2}} = t_{i,j} \cdot g_j^c$ holds, where $r_{i,1}, r_{i,2}$ are the two responses for pair $(i,j)$ and $t_{i,j}$ is the corresponding Okamoto commitment. The correctness of the Okamoto component for inequality relies on the algebraic identity: if $x_i = x_j$, then $(x_i - x_j)^{-1}$ would be undefined (division by zero), making it impossible for the prover to construct valid witnesses. Conversely, when $x_i \neq x_j$, the witnesses $\left(x_i \cdot (x_i - x_j)^{-1}, (x_j - x_i)^{-1}\right)$ satisfy $(g_i \cdot g_j)^{x_i \cdot (x_i - x_j)^{-1}} \cdot (h_i \cdot h_j)^{(x_j - x_i)^{-1}} = g_j$, where we use $h_i = g_i^{x_i}$, $h_j = g_j^{x_j}$, and $(x_j - x_i)^{-1} = -(x_i - x_j)^{-1}$.

Our Rocq formalisation, can be accessed in the file \href{https://anonymous.4open.science/r/SigmaProtocol-EBA5/src/Crypto/NeqSigma.v}{\textcolor{blue}{NeqSigma.v}}, proves the three standard $\Sigma$-protocol properties: (i) \textbf{completeness}:
when the prover knows distinct witnesses $x_1, \ldots, x_{2+n}$ satisfying $g_i^{x_i} = h_i$ and $x_i \neq x_j$
for all $i \neq j$, the verification always succeeds. The proof combines completeness results from both the And composition completeness and the Okamoto protocol completeness for each pair; (ii) \textbf{special honest-verifier zero-knowledge}:
the real transcript distribution (generated with witnesses $x_1, \ldots, x_{2+n}$) is identical to the
simulator distribution (without witnesses). The simulator splits randomness between And and Okamoto
components: for And, it uses the standard simulator producing commitments $t_i = g^{u_i} \cdot h_i^{-c}$;
for each Okamoto pair, it generates simulated commitments independently. Both distributions produce
accepting transcripts with probability $1/|\mathbb{F}|^{(2+n) + (2+n)(1+n)} = 1/|\mathbb{F}|^{(2+n)^2}$, proven
using our uniform probability distribution framework; and (iii) \textbf{special soundness:} from two accepting transcripts sharing identical commitments but responding to distinct challenges \(c_1 \neq c_2\), an extractor is able to reconstruct witnesses \(x_1, \ldots, x_{2+n}\) satisfying the underlying discrete logarithm relations.  However, establishing that these extracted witnesses are genuinely distinct is far more subtle, and cannot, in general, be achieved without additional structural assumptions on the generators. In classical treatments, this is covered by an independence assumption; in our development, we instead expose the underlying algebraic relation explicitly. As formalised in Listing~\ref{neq_special_sound} in Appendix~\ref{appendix_listings}, we obtain an unconditional dichotomy: either the extracted witnesses are pairwise distinct, or the generator structure itself is degenerate in a precise sense, namely that there exists a non-trivial discrete logarithm relation between generators, i.e., \(g_j = g_i^\alpha\) for some \(\alpha\) that is efficiently computable, or one of the generators collapses to the identity.  In this way, special soundness is derived without postulating independence as an axiom. Instead, we exploit decidability of the relevant algebraic predicates to transform what is typically an implicit assumption into an explicit disjunction:
the protocol is sound either because witnesses separate cleanly, or because the algebraic structure itself reveals an underlying dependency. Rocq's insistence on these structural requirements---invisible at the informal level---is precisely what forced this dichotomy into the open.

The implementation \texttt{\seqsplit{generalised\_construct\_neq\_conversations\_real\_transcript}} requires randomness of dimension \((2+n) + (2+n)(1+n) = (2+n)^2\) and produces proof transcripts of size \(O(n^2)\), reflecting the inherently quadratic nature of enforcing pairwise inequality constraints. Consequently, while theoretically sound and fully general, this construction is not particularly practical for large-scale instances, as both the randomness requirements and transcript size grow rapidly with \(n\).

\subsection{Linear Relations}
A further type of subprotocol that we have formalised involves a simple linear relation between committed secret values.
Concretely, we consider the use of Pedersen commitments~\cite{pedersen1991non} for a standard setup given by two generators $g$
and $h$ (where the discrete logarithm of $h$ with respect to $g$ is unknown).

For publicly given Pedersen commitments $C_1, C_2, \ldots, C_n$, publicly given coefficients $\alpha_1, \alpha_2, \ldots, \alpha_n$,
and public value $z$, the prover proves knowledge of values $(v_1, r_1), (v_2, r_2), \ldots, (v_n, r_n)$ such that:
\begin{itemize}
\item $C_i = g^{v_i} \cdot h^{r_i}$ for all $i$ (each commitment opens correctly)
\item $z = \sum_{i=1}^{n} \alpha_i \cdot v_i$ (a linear constraint on the committed values)
\end{itemize}

The protocol follows the standard three-move $\Sigma$-protocol structure for $n$ commitments with a single linear constraint, requiring $2n$ fresh randomness values (denoted $u_1, \ldots, u_n$ and $w_1, \ldots, w_n$) and producing $(n+1)$ commitment group elements and $2n$ response field elements. In the commitment phase, the prover samples fresh randomness $u_1, \ldots, u_n, w_1, \ldots, w_n$ and computes for each $i \in {1, \ldots, n}$ the commitment $t_i = g^{u_i} \cdot h^{w_i}$, mimicking the structure of $C_i$, along with an additional commitment $t_{n+1} = g^{\sum_{i=1}^{n} \alpha_i \cdot u_i}$ to the linear combination of randomness. The prover then sends the commitment vector $(t_1, \ldots, t_n, t_{n+1})$ to the verifier. In the challenge phase, the verifier responds with a random field element $c$. During the response phase, the prover computes for each $i \in \{1, 2, \dots, n\}$ the responses $s_{1,i} = u_i + c \cdot v_i$ and $s_{2,i} = w_i + c \cdot r_i$, sending the two vectors $(s_{1,1}, \ldots, s_{1,n})$ and $(s_{2,1}, \ldots, s_{2,n})$. The verifier accepts the transcript if both verification equations hold: for all $i$, $g^{s_{1,i}} \cdot h^{s_{2,i}} = t_i \cdot C_i^c$ (verifying each Pedersen commitment), and $g^{\sum_{i=1}^{n} \alpha_i \cdot s_{1,i}} = t_{n+1} \cdot g^{c \cdot z}$ (verifying the linear constraint). Verification correctness follows from the homomorphic properties of the commitment scheme: if the
prover is honest with secrets $(v_i, r_i)$ satisfying $C_i = g^{v_i} \cdot h^{r_i}$ and
$z = \sum_i \alpha_i \cdot v_i$, then:
\begin{align*}
g^{s_{1,i}} \cdot h^{s_{2,i}} &= g^{u_i + c \cdot v_i} \cdot h^{w_i + c \cdot r_i} \\
&= (g^{u_i} \cdot h^{w_i}) \cdot (g^{v_i} \cdot h^{r_i})^c = t_i \cdot C_i^c \\
g^{\sum_i \alpha_i \cdot s_{1,i}} &= g^{\sum_i \alpha_i \cdot (u_i + c \cdot v_i)} \\
&= g^{\sum_i \alpha_i \cdot u_i} \cdot g^{c \cdot \sum_i \alpha_i \cdot v_i} = t_{n+1} \cdot g^{c \cdot z}
\end{align*}

Our Rocq formalisation, available at \href{https://anonymous.4open.science/r/SigmaProtocol-EBA5/src/Crypto/PedLinearRel.v}{\textcolor{blue}{PedLinearRel.v}}, implements this protocol with generic support for arbitrary $n \geq 2$ commitments and linear constraints. The formalisation proves all three security properties: (i) \textbf{completeness}: when the prover knows values $(v_1, r_1), \ldots, (v_n, r_n)$ satisfying both $C_i = g^{v_i} \cdot h^{r_i}$ and $z = \sum_i \alpha_i \cdot v_i$, the protocol always verifies successfully; (ii) \textbf{special soundness}: from two accepting transcripts with the same commitments but different challenges $c_1 \neq c_2$,
an extractor recovers the witness values: $v_i = (s_{1,i}^{(1)} - s_{1,i}^{(2)}) \cdot (c_1 - c_2)^{-1}$
and similarly for $r_i$. The extraction is guaranteed to produce valid commitments $C_i = g^{v_i} \cdot h^{r_i}$.
For the linear constraint, the formalisation proves a dichotomy: either $z = \sum_i \alpha_i \cdot v_i$ holds,
or the group element $g$ is the identity (which is a degenerate case). This unconditional result avoids
assumptions on the problem structure; and (iii) \textbf{special honest-verifier zero-knowledge}:
the distribution of real transcripts (generated with honest witnesses) is identically equal to the
simulator distribution (without witnesses). Both produce accepting transcripts with probability
$1/|\mathbb{F}|^{2n}$, proven using our uniform distribution framework. The identical distributions
guarantee information-theoretic security: the verifier learns nothing beyond the validity of the
statements, even with unlimited computational power.

This protocol directly supports confidential transactions~\cite{10.1007/978-3-030-17653-2_22}: committing to an input amount $v_1$, output amount $v_2$, and new balance $v_3$, and proving $v_1 - v_2 - v_3 = 0$ (coefficients $\alpha = (1, -1, -1)$, public value $z = 0$) shows that money is conserved without revealing any amount---the technique used in cryptocurrencies such as Mimblewimble and Monero. We formalise this instance in the \href{https://anonymous.4open.science/r/SigmaProtocol-EBA5/src/Examples/PedLinearRelIns.v}{\textcolor{blue}{PedLinearRelIns.v}} file.

Note that sometimes the direct approach of just revealing $r=\sum_{i=1}^{n} \alpha_i \cdot r_i$ and checking that $g^z h^r = \prod_{i=1}^n C_i^{\alpha_i}$ suffices in standalone protocols (see, e.g, the tallying protocol of~\cite{CFSY96}).
